\scanpage{64}

kum\deleted{é}mérien \inserted{connexe}\begin{tikzpicture}[baseline=-1ex,x=.01cm,y=-.01cm]
\draw[line width=.35pt] (71,0) .. controls (70,21) and (25,13) .. (11,27) .. controls (0,36) and (1,51) .. (3,65);
\end{tikzpicture} (ramifié) de \inserted{$S$, de}\begin{tikzpicture}[baseline=-1ex,x=.01cm,y=-.01cm]
\draw[line width=.35pt] (92,0) .. controls (103,26) and (47,30) .. (28,29) .. controls (10,30) and (0,46) .. (5,64);
\end{tikzpicture} groupe de \deleted{g}Galois $\mu_a$, le diagramme \deleted{étant} induit par le changement de base $\eta\longrightarrow S$ étant cartésien. Un tel diagramme est unique à isomorphisme près, \deleted{$X'$ est connexe,} et $f':X'\longrightarrow S'$ satisfait \margin{\begin{tikzpicture}[x=.01cm,y=-.01cm]
\draw[line width=.35pt] (8,0) -- (2,18);
\draw[line width=.35pt] (0,9) -- (13,15);
\end{tikzpicture}}aux mêmes hypothèses \deleted{ici} que $f$ dans \qquad, mais l'entier $d(X'/S')$ \margin{\begin{tikzpicture}[x=.01cm,y=-.01cm]
\draw[line width=.35pt] (11,0) -- (1,20);
\draw[line width=.35pt] (0,6) -- (8,16);
\end{tikzpicture}}(défini comme $d=d(X/S)$ dans (2)) est maintenant \deleted{une puissance} égal\deleted{e} à $p^r$.
Si $I'\subset I$ est le groupe d'inertie \deleted{de} relatif à $S'$, alors pour tout ent\unsure{i|t}er $n$ premier à $p$, \deleted{\unsure{o|a}}l'homomorphisme
\[
u'_n:H^1(X',\mu_n)\longrightarrow H^1(X'_{\bar\eta},\mu_n)^{I'}
\]
est injectif, et son conoyau est isomorphe à $\deleted{\unsure{\delta}}{}_n\delta/{}_n\delta'$, où $\delta$ est le \margin{\begin{tikzpicture}[x=.01cm,y=-.01cm]
\draw[line width=.35pt] (7,0) -- (0,17);
\draw[line width=.35pt] (0,10) -- (11,17);
\end{tikzpicture}}sous-groupe de torsion de $\operatorname{Coker}w$ ($w$ défini dans (5)) et $\delta'$ le sous-groupe de torsion de $\operatorname{NS}(X_{\unsure{0|o}})$. On en conclut, pour tout nombre premier $\ell\neq p$, \deleted{une} que l'homomorphisme
\[
u'(\ell):H^1(X',\underline Z_\ell\deleted{(\unsure{1|2})})\longrightarrow H^1(X'_{\bar\eta},\underline Z_\ell\deleted{(\unsure{1|2})})^{I'}
\]
est un \uline{isomorphisme}.\makebox[0pt][l]{\begin{tikzpicture}[baseline=-.4ex,x=1mm,y=1mm]\draw[line width=.35pt] (1.1,2.7) .. controls (-.3,1.8) and (-.3,.8) .. (.9,0);\end{tikzpicture}}\deleted{et la même conclusion}

\margin{63}
